\noindent\hspace*{2em}多模态情感识别需要同时利用文本、语音和视觉信息，但三种模态的时间粒度与特征形式不同，原始特征难以直接对应；局部连续缺失会改变可用信息，影响情感分类与强度回归；预测结果还需要说明模型主要依赖哪些模态和局部输入。围绕这些问题，本文依次完成原始视频的三模态特征提取与时序对齐、连续缺失条件下的情感预测，以及模态贡献、模态交互和关键输入位置分析。

针对问题一，我们对附件1的100条原始视频分别提取文本、语音和视觉特征，采用的编码器依次为RoBERTa-base、WavLM-base-plus和SigLIP2-B/16。文本结合语音识别时间信息确定词语或片段的位置，语音由WAV采样范围恢复媒体时间，视觉按采样画面的真实显示时间确定对应区间。文本以词语或片段为单位，语音表示具有较高时间密度，视觉则来自离散采样画面，三种序列在相同索引处未必对应同一段视频内容。因此，时间对应以各观测所覆盖的真实媒体区间为依据，保留各模态原有的时间分布。以媒体时间为共同基准，将每条视频的有效时长划分为50个公共时间窗，按最大时间重叠选择对应的原生特征，同时保存有效掩码和来源记录。映射后各模态具有相同的序列长度，同时通过有效掩码和来源记录保留原始观测的覆盖情况与来源位置。最终得到文本、语音和视觉各$100\times50\times768$的特征数组，平均有效覆盖率分别为0.8242、0.9910和1.0000。典型样本回查表明，公共时间窗能够对应到文本片段、WAV采样范围和视频源帧。

针对问题二，我们使用附件2提供的50位对齐特征，同时预测三分类情感极性和连续情感强度。掩码均值池化模型MMP在有效位置上等权聚合，各模态表示经独立投影后拼接，分别进入分类头与回归头，作为分类与回归基线；考虑到不同位置的情感信息并不等量，在其上加入内容驱动的时间加权残差分支，得到轻量时间注意力残差池化模型LTARP。均值分支保留有效序列的整体信息，加权分支补充不同位置的内容差异，可学习残差系数控制补充表示的作用。分类与回归采用加权交叉熵和SmoothL1损失联合优化，使同一融合表示同时服务于极性判断和强度估计。该模型仅增加883个参数，总参数量为164343。从缺失模态、缺失位置和缺失比例三个方面构造54种连续局部缺失场景，比较完整输入与遮挡输入下的分类和回归结果。场景覆盖单模态及双模态缺失，分别考察连续信息损失发生在前段、中段或后段时的影响。评价同时采用宏平均F1、平均绝对误差MAE及Pearson相关系数，分别反映类别识别、连续强度误差和预测变化的一致程度。完整输入下，LTARP的宏平均F1、MAE和Pearson分别为0.6243、0.6052和0.6475；54种缺失场景平均下分别为0.6163、0.6086和0.6417。分模态结果显示，文本连续缺失带来的退化最明显，语音和视觉缺失的影响相对较小，表明三种信息在该组样本中的预测作用存在差异。为考察LTARP与不同多模态建模路线的相对表现，将11种公开方法适配到相同对齐特征和双任务输出接口。在本文统一输入与评价设置下，LTARP在约16.4万参数规模下保持了较高的宏平均F1，并兼顾连续强度预测和连续缺失场景表现。固定模型完成附件3全部30条无标签样本的情感极性与连续强度预测，输出对应的逐样本结果。

针对问题三，我们固定问题二的预测模型，枚举三种模态的全部8种输入组合并计算精确Shapley贡献，分别分析文本、语音和视觉对分类及回归输出的作用。分类分析围绕固定预测类别的对数优势展开，回归分析则围绕连续预测输出展开，从而分别描述各模态对两类任务的作用。进一步比较两两模态联合输入的贡献与简单可加组合，考察不同信息组合对当前预测目标的影响。交互正值表示联合输入的贡献高于可加组合，负值表示低于可加组合，以此区分单一模态作用与联合输入产生的额外变化。模态贡献和交互描述整路输入的作用，为进一步确定重要信息位于序列中的具体位置，在分类主导模态内采用连续窗口遮挡，保持其余输入不变，逐一比较候选区间被遮挡前后的输出，定位高影响位置。以随机等长窗口作为对照，使两类区间的输出变化在相同遮挡长度下进行比较；再通过删除实验观察逐步移除所选位置后的响应，检查关键位置分析与预测变化的一致性。独立验证子集中，遮挡所选窗口后的分类对数优势平均下降0.482771，随机等长窗口下降0.082036，差值为0.400735，按视频分组自助采样得到的95\%区间为$[0.3635,\,0.4372]$。所选区间被遮挡后对分类输出的影响明显大于随机等长区间。解释结果保留模态贡献、关键区间和对应输入来源：文本关键位置可以回查到原文字符区间，语音和视觉关键位置可以回查到相应未对齐特征行。最终完成附件4全部20条无标签样本的预测与解释，逐样本报告情感极性、连续强度、三模态贡献和关键输入区间。

综上，三模态特征能够在统一时间位置上保留原始来源对应关系，LTARP以较小参数增量完成局部连续缺失条件下的情感分类与强度回归，Shapley贡献和连续遮挡进一步确定预测依赖的模态与关键输入位置。所得到的结果覆盖原始特征构建、无标签样本预测和输入作用分析，并保留了可供逐样本核对的来源信息，为复杂场景下多模态情感识别提供了具体方法及可复现结果。

\par\vspace{0.7\baselineskip}
\noindent\textbf{关键词：}多模态情感识别；时序对齐；局部连续缺失；情感预测；Shapley值；可解释分析
